% Editable, self-contained LaTeX source. Compile with pdfLaTeX (twice) or latexmk.
\documentclass[10pt,a4paper]{article}
\usepackage[T1]{fontenc}
\usepackage{lmodern}
\usepackage[a4paper,left=25mm,right=25mm,top=24mm,bottom=24mm]{geometry}
\usepackage{amsmath,amssymb,array,booktabs,xcolor,microtype}
\usepackage[labelfont=bf,font=small]{caption}
\usepackage{cite}
\usepackage{hyperref}

\definecolor{citeblue}{HTML}{1B4068}
\definecolor{diaggray}{HTML}{5A636A}

\hypersetup{
  colorlinks=true,
  linkcolor=citeblue,
  citecolor=citeblue,
  urlcolor=citeblue,
  pdftitle={Focusing Clocks, Closed-Trace Sewing, and State Counts in Null-Cone Quantum Gravity},
  pdfauthor={Null-Cone Gravity Study},
  pdfsubject={Analysis of arXiv:2611.34924v1}
}

\linespread{1.05}
\numberwithin{equation}{section}
\allowdisplaybreaks[2]
\clubpenalty=2500
\widowpenalty=2500
\displaywidowpenalty=2500

\newcommand{\repo}[2]{\href{https://github.com/sprajs/quantum-gravity/blob/main/#1}{#2}}
\newcommand{\tr}{\operatorname{Tr}}
\newcommand{\rank}{\operatorname{rank}}
\newcommand{\HH}{\mathcal{H}}

\title{\Large\bfseries Focusing Clocks, Closed-Trace Sewing, and State Counts\\in Null-Cone Quantum Gravity}
\author{Null-Cone Gravity Study\thanks{AI-assisted theoretical study of V.~Chandrasekaran, \href{https://arxiv.org/abs/2611.34924v1}{\emph{An algebraic quantization of gravity from the null cone}}, arXiv:2611.34924v1 (25~November~2026). See Section~\ref{sec:repro} for authorship, independence, and reproducibility disclosures.}}
\date{26 November 2026}

\begin{document}
\maketitle

\begin{abstract}
Does a bound on the size of a causal diamond limit the number of quantum states it can contain? We examine this question in Chandrasekaran's null-cone quantization of four-dimensional vacuum gravity (arXiv:2611.34924v1). Classical focusing bounds the corner area in terms of an affine clock. The affine bulk curvature density reproduces the Einstein action, and the single-harmonic Gaussian graviton-pair overlap sums to a dilogarithm. At finite cap size the proposed symplectic form has a nonzero exterior derivative. Varying the cap location recovers the missing term on the fixed-label shear sector. In a flat diamond, two narrowing vacuum quadrupoles on opposite sides of the cap rim leave its local geometry and that of the waist unchanged to first order. Their joined Einstein pairing with every fixed smooth perturbation tends to zero, while the mixed variation of the clock retains a finite contribution from the outer packet. This excludes a continuously differentiable family of smooth Hamiltonian vectors through flat space for the clock excess with local rim charges. Varying the radial profiles extends the obstruction to boundary terms with finitely many independent responses in this angular harmonic. If a uniform exterior boost generates the waist area, the same obstruction extends to the clock itself with freely varying area. In the quantum state count, the continuum clock must first admit a state below the prescribed reading. Even if it does, the reflected weight has a strictly positive lower bound on the accepted space, so finiteness of its ordinary trace is equivalent to finite rank. Convergence of amplitudes between normalizable states leaves that trace undetermined. A two-area model and a finite two-copy calculation isolate this trace problem. We derive the assumptions under which the Schwarzschild area packet reproduces Sen's spin-summed logarithmic coefficient. In a matter extension, reversing a scalar vacuum--pair phase reverses the leading semiclassical clock excess on both sheets at fixed mean energy.
\end{abstract}

\section{Introduction}
\label{sec:intro}

Two events on an observer's worldline select a finite spacetime region: light emitted from the past event can scatter inside and return by the future event.
The domain accessible between these events is a \emph{causal diamond}, bounded by an expanding null cone from the past tip and a contracting null cone toward the future tip that intersect on a spacelike two-sphere, the corner (Fig.~\ref{fig:diamond}).
In vacuum general relativity, gravitational shear distorts a bundle of neighbouring null geodesics and slows its expansion, so rays travel farther in affine parameter to open out to a prescribed corner area~\cite{Raychaudhuri1955,Sachs1961}.
In Ref.~\cite{Chandrasekaran2026} (version~1, 83~pages, 25~November~2026), Chandrasekaran turns this geometric run from a finite cap to the corner into an intrinsic gravitational clock and proposes an algebraic, nonperturbative quantization of four-dimensional vacuum gravity from data on the intersecting null cones.

The central physical question of Ref.~\cite{Chandrasekaran2026} is how many \emph{perfectly distinguishable} preparations---represented by mutually orthogonal quantum states---can fit inside a causal diamond whose geometric clock reading is bounded above.
Classically, an upper bound on the clock reading enforces an upper bound on the corner area.
In the quantum theory, counting distinguishable preparations requires two independent steps: showing that the quantum clock admits physical states below the specified threshold, and bounding the number of linearly independent orthogonal states that share an allowed reading.

\begin{figure}[tbp]
\centering
\setlength{\unitlength}{0.82mm}
\begin{picture}(161,64)
\color{black}\linethickness{1.0pt}
\put(51,4){\line(1,1){27}}\put(51,4){\line(-1,1){27}}
\put(24,31){\line(1,1){27}}\put(78,31){\line(-1,1){27}}
\color{diaggray}\linethickness{0.45pt}
\multiput(25,31)(3,0){18}{\line(1,0){1.5}}
\put(45,10){\line(1,0){12}}\put(45,52){\line(1,0){12}}
\put(12,9){\vector(0,1){44}}
\put(4,55){\makebox(16,4){\small time}}
\color{black}\put(51,4){\circle*{1.8}}\put(51,58){\circle*{1.8}}
\put(78,31){\circle*{1.8}}\put(24,31){\circle*{1.8}}
\color{diaggray}\put(55,58){\line(1,0){26}}
\put(81,31){\line(1,0){7}}\put(55,4){\line(1,0){26}}
\put(84,58){\makebox(0,0)[l]{\small Future tip (final reception)}}
\put(91,31){\makebox(0,0)[l]{\small Corner where the cones meet}}
\put(84,4){\makebox(0,0)[l]{\small Past tip (initial flash)}}
\end{picture}
\caption{Causal diamond bounded by intersecting past and future null cones. A spatial slice runs horizontally; the two outer vertices at the widest slice represent a two-sphere in three spatial dimensions, and the short horizontal segments near the past and future tips mark the finite spacelike caps.}
\label{fig:diamond}
\end{figure}

The finite-cap obstruction below concerns the symplectic form of version~1 on its stated phase space. The quantum examples have a different role. They test the inferences from a bounded clock to finite rank and from sewing between two given states to a finite trace. Their conclusions about the full gravitational space depend on whether its constrained observables realize the hypotheses of the models.

Unless stated otherwise, we adopt natural units with the speed of light $c=1$ and Boltzmann's constant $k_B=1$, writing the squared Planck length as $\ell_P^2=G_N\hbar$ and the dimensionless corner area as $s=A/(4G_N\hbar)$, where $G_N$ is Newton's constant and $\hbar$ is the reduced Planck constant.
SI units (metres, seconds, kilograms, hertz) are restored explicitly in Section~\ref{sec:observables} when evaluating operational benchmarks.
Numbered displays such as~\eqref{eq:focusing} refer to equations in the present paper; equations in the studied manuscript are cited explicitly by section or equation number in Ref.~\cite{Chandrasekaran2026}.

Sections~\ref{sec:classical}--\ref{sec:clock} establish the geometric and spectral facts needed for the trace argument in Section~\ref{sec:sewing}. We then examine its consequences for black-hole entropy and for measurements with a specified preparation.

\section{Classical Null-Cone Focusing and the Finite-Cap Phase Space}
\label{sec:classical}

\subsection{Focusing equation and the affine clock}

Label the null generators of an expanding light cone by their launch directions on a fiducial unit sphere.
On a small bundle of rays, write the cross-sectional area per unit solid angle as $r^2$, so that $r(u)$ is its area radius along an affine parameter $u$ for which freely moving null rays obey the geodesic equation without parameter-dependent acceleration, and fix a boundary frame to compare directions and readings~\cite{Reisenberger2008,Reisenberger2013,ChandrasekaranFlanagan2023}.
On a twist-free bundle in four-dimensional vacuum Einstein gravity (with or without a cosmological constant), the Raychaudhuri focusing equation~\cite{Raychaudhuri1955,Sachs1961} reduces to
\begin{equation}
  \frac{d^2 r}{du^2} = -\frac{r}{2}\,\sigma_{AB}\sigma^{AB}.
  \label{eq:focusing}
\end{equation}
The traceless shear $\sigma_{AB}$ measures the rate of shape distortion of the bundle at fixed cross-sectional area.
Because $\sigma_{AB}\sigma^{AB}\ge 0$, gravitational radiation that induces shear reduces the expansion rate $r'(u)$, while a round shear-free bundle satisfies $r''=0$.

Following Section~3.2 of Ref.~\cite{Chandrasekaran2026}, truncate the expanding cone from the past tip (the manuscript's incoming cone) at a spacelike cap of area radius $\epsilon r_S$ on each ray, where $0<\epsilon<1$ and $r_S$ is the corner radius.
Let $u_c$ and $u_S$ denote the affine coordinates at the cap and the corner, respectively.
Integrating~\eqref{eq:focusing} twice yields
\begin{equation}
  u_* \equiv r'(u_c)(u_S - u_c)
  = (1-\epsilon)r_S
  + \frac{1}{2}\int_{u_c}^{u_S} (u_S - w)\,r(w)\,\sigma_{AB}\sigma^{AB}(w)\,dw.
  \label{eq:affine-clock}
\end{equation}
Rescaling the affine parameter $u\mapsto a u + b$ ($a>0$) rescales the cap slope $r'(u_c)$ and the interval $u_S-u_c$ inversely, leaving $u_*$ invariant; the integral supplies the strictly nonnegative excess reading caused by shear.
For the angularly averaged clock $\bar{u}_*$ weighted by $r_S$ over a corner of total area $A$ (Ref.~\cite{Chandrasekaran2026}, Eq.~(3.28)), Cauchy--Schwarz gives the classical area ceiling
\begin{equation}
  \bar{u}_* \ge (1-\epsilon)\sqrt{\frac{A}{4\pi}},
  \label{eq:clock-bound}
\end{equation}
which is saturated precisely by a round shear-free sphere.

Two geometric points qualify the scope of~\eqref{eq:affine-clock}--\eqref{eq:clock-bound} (\repo{audit/03-2.md}{focusing derivation}).
First, on any smooth positive-area bundle the integrated identity~\eqref{eq:affine-clock} and inequality~\eqref{eq:clock-bound} survive even if $r(u)$ passes through a local maximum ($r'=0$) prior to a caustic.
However, using the dimensionless area ratio $v=r/r_S$ as a global coordinate along the entire ray---as in the single-chart phase-space reduction of Ref.~\cite{Chandrasekaran2026}, Section~3.2---requires strict monotonicity $r'>0$; an explicit vacuum plane wave reaches $r'=0$ before its first caustic and makes the inverse $u(v)$ double-valued.
Second, because the focusing integral depends only on radial derivatives of the conformal two-metric along each ray, Eq.~\eqref{eq:clock-bound} imposes no classical restraint on arbitrarily rapid transverse angular oscillations of the shear.

\subsection{Bulk equations and the finite-cap symplectic obstruction}

On the branch with an invertible gravitational tetrad, the constrained first-order bulk equations of Ref.~\cite{Chandrasekaran2026}, Section~3.1, reduce to vacuum Einstein's equations, and the smooth-vertex corner potential (Eq.~(3.42a) of Ref.~\cite{Chandrasekaran2026}) yields the expected corner boost charge $A/(8\pi G_N)$~\cite{ChandrasekaranFlanagan2023}.
Introducing a finite cap at fixed area ratio $v_c=\epsilon$, however, exposes a concrete \emph{defect in the construction as written}.

To define Poisson brackets obeying the Jacobi identity, the reduced symplectic two-form $\omega$ must be closed ($d\omega=0$) on its declared space of variations.
Consider variations allowed by the fixed-ratio cap condition (Ref.~\cite{Chandrasekaran2026}, Eq.~(3.18)): vary the corner area density together with two independent radial profiles of a single shear polarization while holding angular maps fixed.
Taking the field-space exterior derivative of the displayed double-null two-form (Ref.~\cite{Chandrasekaran2026}, Eq.~(A.14)) leaves an uncancelled three-form localized at the cap rim $v=\epsilon$ (\repo{audit/03-3.md}{closure test}; \repo{audit/cap-crosscheck.md}{cap-variation check}).
Let $a=r_S^2$ denote the local corner area density per unit fiducial solid angle. In the single-polarization sector with cap fraction $b=\epsilon$, the rim remainder equals $\tfrac{1}{2}b^3\,da\wedge dz'(b)\wedge dz(b)$ in the normalization with $(16\pi G_N)^{-1}$ and the fiducial angular measure factored out.
Evaluated on the three-parameter polynomial family $z(v)=x(1-v)^2+y(1-v)^3$ (whose corner value and first derivative vanish at $v=1$), the exterior derivative on $(\partial_a,\partial_x,\partial_y)$ equals $b^3(1-b)^4/2$, which vanishes at the smooth vertex ($b=0$) and equals $1/256$ at $b=1/2$, as confirmed in exact rational arithmetic by \repo{checks/03-3.py}{\texttt{checks/03-3.py}}.
By contrast, Reisenberger's closed double-null two-form~\cite{Reisenberger2008,Reisenberger2013} holds the \emph{absolute} area density fixed on each truncation cut rather than the ratio $b=\epsilon$.
Neither adding the exterior derivative of the cap-joint potential (Ref.~\cite{Chandrasekaran2026}, Eq.~(A.18)) nor vacuum evolution linking past and future caps cancels this rim term.

Any covariant repair that retains moving caps and their joints, or narrows the admissible variations, must also produce a consistent Hamiltonian generator (Ref.~\cite{Chandrasekaran2026}, Section~3.4 and Appendix~B).
While the bulk first-order Hamiltonian reduces on the constraints to a boundary integral and a fixed-location corner boost integrates cleanly to $A/(8\pi G_N)$, a translating cut carries an additional boundary flux $-i_\xi\theta$ in the covariant charge variation; on a Schwarzschild sphere the surface integral in the stated normalization equals $M/2$, giving $\delta M/2$ in its variation, and on unrestricted radiating null data the fixed-embedding bulk charge variation is nonintegrable without compensating boundary degrees of freedom or flux-restricting boundary conditions (\repo{audit/03-4.md}{Hamiltonian and boundary-flux analysis}).

On the torsion-free metric branch the full affine identity (B.10)--(B.11) also holds off shell for arbitrary smooth angular metric, shift and $F$. Splitting the four-dimensional curvature into the intrinsic cut curvature, expansion, shear, twist and divergence recovers the three printed boundary currents. An independent symbolic contraction matches all 39 monomials after normalizing the angular frame at one point while retaining arbitrary first and second derivatives of the metric functions (\repo{audit/bulk-identity.md}{general proof and exact symbolic check}). In (B.7) the angular divergence integrates to zero on a closed cut while the radial derivative contributes at its endpoints.

\section{Single-Cone Graviton-Pair Overlap and Weak Coupling}
\label{sec:pair}

Section~4.4 of Ref.~\cite{Chandrasekaran2026} (Eqs.~(4.32)--(4.34)) furnishes a quantum calculation on a single incoming null cone that closes in exact analytic form.
Restricting to a single angular pattern of gravitational radiation and one polarization, radial oscillations carry integer mode labels $k=1,2,\dots$ represented by independent Gaussian oscillators in a product vacuum $\Omega$.
With $c=1$, define
\begin{equation}
  s = \frac{A}{4G_N\hbar}, \qquad x = \frac{1}{s}, \qquad v_k = \langle q_k^2\rangle = \frac{x}{2k}.
  \label{eq:pair-params}
\end{equation}
At fixed corner area $A$, decreasing $x$ weakens the coupling.
The vacuum-subtracted source $q_k^2 - v_k$ acts on $\Omega$ to create a pair of quanta $|2_k\rangle$ in mode $k$ with amplitude $\sqrt{2}\,v_k$, and distinct pair states are mutually orthogonal.
Applying the chosen coupling-dependent affine displacement multiplies each $|2_k\rangle$ by the phase $e^{-2ikx}$.

\subsection{Exact mode summation}

With radial modes retained through $k=N$, the overlap of the pair preparation with its displaced copy is
\begin{equation}
  C_N(x) = 2\sum_{k=1}^{N} v_k^2\, e^{-2ikx}
  = \frac{x^2}{2}\sum_{k=1}^{N} \frac{e^{-2ikx}}{k^2}.
  \label{eq:pair-sum}
\end{equation}
Every summand has modulus $x^2/(2k^2)$, so the continuum limit $C(x)=\lim_{N\to\infty}C_N(x)$ converges absolutely for all real $x$, with universal tail bound $|C(x)-C_N(x)|\le x^2/(2N)$ (and a sharper $O(N^{-2})$ oscillatory bound away from $x\in\pi\mathbb{Z}$).
In closed form,
\begin{equation}
  C(x) = \frac{1}{2}\,x^2 \operatorname{Li}_2\!\left(e^{-2ix}\right),
  \qquad
  \operatorname{Li}_2(z) = \sum_{k=1}^{\infty}\frac{z^k}{k^2}\quad (|z|\le 1).
  \label{eq:pair-limit}
\end{equation}
At $x=0.1$, Eq.~\eqref{eq:pair-limit} evaluates to $C(0.1) \simeq 0.006703874 - 0.002609994\,i$; a direct sum of~\eqref{eq:pair-sum} through $N=100{,}000$ radial modes agrees with an independent quadrature evaluation to within $2.50\times 10^{-12}$, inside the analytic tail bound (\repo{audit/04-pair.md}{derivation and error bounds}; \repo{checks/pair.py}{\texttt{checks/pair.py}}).

\subsection{Nonuniform weak-coupling expansion and physical scope}

Expanding each phase as $e^{-2ikx} = 1 - 2ikx + O(k^2 x^2)$ at fixed cutoff $N$ yields an $x^3$ coefficient of $-iH_N$, where $H_N = \sum_{k=1}^{N} k^{-1} = \log N + \gamma + O(N^{-1})$ diverges without bound as $N\to\infty$.
In the exact continuum sum~\eqref{eq:pair-limit}, differentiating $\operatorname{Li}_2(e^{-i\theta})$ along the unit circle gives, as $x\downarrow 0$,
\begin{equation}
  \operatorname{Im} C(x) = x^3\bigl[\log(2x) - 1\bigr] + O(x^5),
  \label{eq:pair-asymp}
\end{equation}
accompanied by a one-sided real cubic term $-\pi x^3/2$ that vanishes identically at every finite $N$.
Both iterated limits of the overlap itself are zero. Extracting its cubic Taylor coefficient at finite $N$ fails to commute with removing the cutoff because modes with $kx\sim 1$ accumulate order-one phases.

Although this derivation reproduces Eqs.~(4.32)--(4.34) of Ref.~\cite{Chandrasekaran2026} completely, its physical reach is narrower than an interacting gravitational scattering amplitude (\repo{audit/scattering.md}{scattering comparison}).
The overlap~\eqref{eq:pair-sum} is already finite at zero displacement ($C=\pi^2 x^2/12$ at zero phase) because the Gaussian pair weights fall as $k^{-2}$, and the $x^3\log x$ nonanalyticity enters only because the chosen affine displacement shrinks linearly with $x$---holding the displacement fixed independently of $x$ leaves a constant times $x^2$ at fixed area.
Because the example uses independent free oscillators on a single incoming cone without Einstein interaction vertices or outgoing wave packets, testing genuine gravitational self-interactions---such as reproducing the universal two-loop four-graviton all-plus scale logarithm associated with the Goroff--Sagnotti $R^3$ divergence~\cite{GoroffSagnotti1986,DunbarJehuPerkins2017}---requires constructing interacting evolution on external wave packets in a much larger region.

\section{Quantum Clock Spectrum, Shape Modes, and Ultraviolet Shear}
\label{sec:clock}

Suppose the continuum quantum clock admits a positive self-adjoint realization $T$ (Ref.~\cite{Chandrasekaran2026}, Sections~4.3 and~5.1).
A maximum allowed reading $\ell$ selects the spectral subspace onto which $P_\ell = \mathbf{1}_{[0,\ell]}(T)$ projects.
For a normalized preparation $\Omega$, two distinct quantities enter the state count: the passing probability $\|P_\ell\Omega\|^2$, which tests whether the clock window is nonempty on $\Omega$, and the rank $\rank(P_\ell)$, which counts mutually orthogonal allowed states.

Section~2 of Ref.~\cite{Chandrasekaran2026} establishes a clean control in two-dimensional JT gravity coupled to a free scalar field (where gravitational shear is absent): fixing both corner dilaton charges and the matter state still leaves a continuous clock coordinate with an unconstrained conjugate momentum, admitting infinitely many mutually orthogonal states with the exact same bounded clock probability distribution and yielding a type~$\mathrm{I}_\infty$ regional algebra (\repo{audit/02-jt.md}{JT gravity analysis}).
In four dimensions, shear enters the clock~\eqref{eq:affine-clock}, and explicit models separate the roles of corner shape degrees of freedom and short-distance shear fluctuations.

\subsection{Shear-free corner-shape benchmark}

To test how corner shape behaves at fixed total area, fix $A=4\pi b^2$ and set the shear to zero (\repo{audit/regulator-program.md}{exact shape benchmark}).
Let a real parameter $q\in\mathbb{R}$ redistribute area over the corner sphere with density proportional to $e^{q\cos\theta}$ ($\theta$ being the polar angle), normalized to keep the total area $A$ fixed, and take wavefunctions in $L^2(\mathbb{R},dq)$.
With $u_0=(1-\epsilon)b$, performing the angular integrals in the clock expression of Ref.~\cite{Chandrasekaran2026} (Eq.~(4.26)) produces the exact multiplication operator
\begin{equation}
  \frac{T(q)}{u_0} = \sqrt{\frac{q}{2}\coth\!\left(\frac{q}{2}\right)},
  \qquad T(0) = u_0.
  \label{eq:shape-clock}
\end{equation}
At ceiling $\ell = 1.1\,u_0$, the accepted region is the interval $|q| \le q_* \simeq 1.6212273$.
For a normalized Gaussian state with probability density $e^{-q^2}/\sqrt{\pi}$, the passing probability is $\|P_\ell\Omega\|^2 = \operatorname{erf}(q_*) \simeq 0.97814$.
Dividing $[-q_*,q_*]$ into arbitrarily many disjoint subintervals and taking their normalized indicator functions yields arbitrarily many mutually orthogonal allowed states, so the exact continuum rank is infinite.
In finite $n\times n$ Hermite-matrix approximations to $q$, the accepted eigenvalue count roughly doubles each time $n$ quadruples (Table~\ref{tab:corner-clock}).

\begin{table}[htbp]
\centering
\caption{Accepted eigenvalue counts for the shear-free corner-shape clock~\eqref{eq:shape-clock} at threshold $\ell = 1.1\,u_0$ ($|q|\le q_*\simeq 1.6212273$) using $n\times n$ Hermite-basis truncations of $q$ (\repo{checks/corner_clock.py}{\texttt{checks/corner\_clock.py}}). The exact continuum rank is infinite.}
\label{tab:corner-clock}
\begin{tabular}{@{}lrrrrr@{}}
\toprule
Truncated matrix dimension $n$ & 16 & 64 & 256 & 1,024 & 4,096 \\
\midrule
Accepted eigenvalues ($\le 1.1\,u_0$) & 6 & 12 & 24 & 46 & 94 \\
\bottomrule
\end{tabular}
\end{table}

This model illustrates why both the Hilbert-space measure and the physical meaning of $q$ matter.
At $\ell=u_0$ only the single point $q=0$ passes, which has zero measure, whereas raising the threshold by $10\%$ admits an entire interval.
Classically, in the smooth-vertex form of Ref.~\cite{Chandrasekaran2026} (Eq.~(3.42a)), a nonconstant smeared local corner area density generates a shift of the relative boost angle alone and leaves the clock-bounded region invariant.
Transferring this infinite-rank mechanism to physical four-dimensional quantum gravity requires three conditions (\repo{audit/finite-kinematics.md}{physical corner-charge tests}): exhibiting a physical shape observable that survives the gravitational constraints and gauge quotient, establishing a continuous spectral component on an allowed state, and proving that its spectral projections commute strongly with the physical clock $T$.
Whether the unspecified canonical chart of Ref.~\cite{Chandrasekaran2026} satisfies these conditions remains open.

\subsection{Ultraviolet shear fluctuations and existence of clock states}

On an expanding light cone, the clock integrates the exponential of accumulated squared shear along each generator.
Squaring a fluctuating field at a point is ultraviolet singular, and replacing a quantum resolvent expectation by a classical Gaussian phase-space tail (Ref.~\cite{Chandrasekaran2026}, Eq.~(5.16), noted as heuristic in its Footnote~18) does not control noncommuting operators in the continuum limit.
If the cutoff clocks define an increasing family of closed positive forms on a common Hilbert space, the monotone convergence theorem~\cite{Behrndt2010} gives a limiting self-adjoint clock $T$ on the closure $\mathcal{M}=\overline{D(Q)}$ of its finite-form domain, with $\lim_{\ell\to\infty}\|P_\ell\Omega\|^2=\|P_{\mathcal{M}}\Omega\|^2$, but does not rule out $\mathcal{M}=\{0\}$.

Two explicit model checks sharpen this ultraviolet obstruction (\repo{audit/shear-uv.md}{short-distance shear analysis}).
First, given the linearized null-plane shear commutator derived in Appendix~D.2 of Ref.~\cite{Chandrasekaran2026}, no nonzero vector admits a finite bare point-field squared-shear integral over any open radial interval.
Second, with selected bare Gaussian radial modes matching the paper's short-wave symplectic scaling, the vacuum mean squared radial derivative grows as $O(N^2)$ with radial cutoff $N$.
While a diverging vacuum mean alone leaves the probability of a low reading undetermined, in two regulated quadratic shear models even the state-independent \emph{lowest clock eigenvalue} diverges as the radial cutoff and the focusing cap are both removed, emptying any fixed window $P_\ell$ identically; whether the full nonlinear two-polarization gravitational clock behaves similarly remains open.
Vacuum subtraction alters the positivity argument behind~\eqref{eq:clock-bound} because local Wick-subtracted squares can have negative expectations in squeezed states.
A constructive proof of state existence must therefore exhibit a normalized physical state $\chi$ with finite limiting clock form energy $Q[\chi]<\infty$, control its overlap $|\langle\Omega,\chi\rangle|>0$ with the prepared state, and bound $\|P_\ell\Omega\|^2$ from below at the specified $\ell$.

\section{Open Sewing, Closed Traces, and Boost Contrasts}
\label{sec:sewing}

\subsection{Open sewing versus trace closure}

Cutting an open transition amplitude between specified initial and final normalizable preparations and inserting a complete orthonormal basis $\{|j\rangle\}$ on the intermediate cut yields a sum of products $a_j^* b_j$, where $a_j$ and $b_j$ are the components of the two finite-norm states at the cut.
By the Cauchy--Schwarz inequality,
\begin{equation}
  \sum_j |a_j^* b_j|
  \le \left(\sum_j |a_j|^2\right)^{1/2}
      \left(\sum_j |b_j|^2\right)^{1/2}
  < \infty.
  \label{eq:cauchy-schwarz}
\end{equation}
This \emph{open sewing} converges absolutely for any unitary evolution even in an infinite-dimensional Hilbert space, for one intermediate cut. For finitely many cuts, successive strong limits of the cut projections recover the composed amplitude. Joint absolute convergence of the multiple series requires further hypotheses (\repo{audit/04-gluing.md}{open-gluing analysis}).
However, absolute convergence in~\eqref{eq:cauchy-schwarz} depends critically on the fixed finite-norm endpoints.
Composing two identity operators between fixed unit vectors yields their finite inner product, whereas closing the same identity operator into a loop demands $\tr(1)=\sum_j 1$, which diverges in infinite dimension.

\subsection{Reflected weights and equivalence of finite trace to finite rank}

Assume now that the four-dimensional quantum clock $T$ exists, commutes strongly with the dimensionless corner area operator $\hat{s}=\hat{A}/(4G_N\hbar)$, and enforces the classical area ceiling~\eqref{eq:clock-bound} on $P_\ell\HH$, so that $0\le \hat{s}\le s_*$ with
\begin{equation}
  s_* = \frac{\pi\ell^2}{(1-\epsilon)^2 G_N\hbar}.
  \label{eq:s-max}
\end{equation}
Centre a Gaussian area filter at $s_b = \pi b^2/(G_N\hbar)$, the value for a round sphere of radius $b$.
Any fixed positive width suffices here; the entropy calculation in Section~5.3 of Ref.~\cite{Chandrasekaran2026} adopts variance $2s_b$, giving the reflected weight
\begin{equation}
  W = e^{-\hat{s}} f_b(\hat{s}), \qquad
  f_b(s) = \exp\!\left[-\frac{(s-s_b)^2}{4s_b}\right].
  \label{eq:reflected-weight}
\end{equation}
Because $s\mapsto e^{-s}f_b(s)$ is continuous and strictly positive on the compact interval $[0,s_*]$, its minimum $c = \min_{s\in[0,s_*]} e^{-s}f_b(s)$ is strictly positive while $W\le 1$.
For every finite-rank projection $F$ within the clock-selected space $P_\ell\HH$, it follows that
\begin{equation}
  c\,\rank F \le \tr(FWF) \le \rank F.
  \label{eq:trace-rank}
\end{equation}
Every independent allowed state adds at least $c>0$ to the ordinary unnormalized trace.
Under these hypotheses, finiteness of the ordinary closed trace $\tr(W)$ is already equivalent to finite rank of $P_\ell$, so an open-sewing bound such as~\eqref{eq:cauchy-schwarz} cannot supply the missing step; independent control of the closed experiment is required.

\subsection{Order dependence in the boost-contrast test}

Section~5.2 of Ref.~\cite{Chandrasekaran2026} attempts to force finite rank by contradiction using a relative boost applied to one side of a reflected cut.
The smooth linearized vacuum has a nonzero response to a relative boost. The contradiction requires that response to survive in the full uncut preparation.
Under the hypothesis $\dim(P_\ell\HH)=\infty$, a finite partition of $[0,s_*]$ has a narrow area cell whose spectral projection has infinite rank, allowing basis states to be grouped into long cycles dominated by nearly equal areas; on each cycle projection $P$, the boost supplies almost a common phase to the dominant states, suppressing the boost contrast and reversing the sign of an unnormalized subtracted observable $X$ relative to the uncut answer (Ref.~\cite{Chandrasekaran2026}, Eq.~(5.38)).

Both the finite-projection cycle bounds and the freely falling vacuum boost signal hold under their stated premises (\repo{audit/dimension-crosscheck.md}{boost-argument cross-check}).
The unproved step is a closed-sewing theorem showing that \emph{every} exhaustive cut order converges to the same independently defined uncut closed amplitude.
For sufficiently small cycle tolerance in Eq.~(D.25) of Ref.~\cite{Chandrasekaran2026}, the estimated normalized signed insertion satisfies $Z_P[X]/Z_P\leq-\kappa$ for a fixed $\kappa>0$, where $Z_P=\tr(PWP)$. Then $|Z_P[X]|\geq\kappa c\,\rank P$ by~\eqref{eq:trace-rank}. Bounded signed sums along this exhaustive family already imply finite rank, independently of the sign of the uncut amplitude (\repo{reviews/fable-clock-sewing.md}{cycle estimate}).

An explicit comparison space isolates this operator-limit dependence (\repo{checks/05-2.py}{\texttt{checks/05-2.py}}).
Take $\HH = \mathbb{C}^2\otimes \ell^2(\mathbb{N})$ with two area eigenvalues $s_0,s_1$, each of infinite multiplicity, positive weights $w_i=e^{-s_i}f_b(s_i)$, and a bounded two-sided exchange observable. The example uses an independent unit-variance Gaussian centered at $2+\pi$, for which the reflected weights coincide at $s_0=1$ and $s_1=1+2\pi$.
This model obeys exact open sewing~\eqref{eq:cauchy-schwarz} for all fixed finite-norm vectors.
Yet three distinct ways of exhausting the same orthonormal basis---varying the ratio of states drawn from the two area sectors ($k=1,4,40$)---give normalized boost contrasts of approximately $-0.460$, $-0.184$, and $-0.022$ ($-0.45970$, $-0.18388$, and $-0.02242$).
Because strong convergence $P\to 1$ does not select a unique limit for an infinite-trace closed contraction, gravitational dynamics would have to supply an additional closed-sewing identity fixing the physical result.

\section{Black-Hole Entropy, One-Loop Coefficients, and Ensembles}
\label{sec:entropy}

In gravitational units ($c=k_B=1$), the Bekenstein--Hawking entropy of a large black hole is its horizon area in Planck units, $s=A/(4G_N\hbar)$~\cite{Bekenstein1973,Hawking1975}.
In Section~5.3 of Ref.~\cite{Chandrasekaran2026}, analytically continuing a half-turn generated by the corner boost charge $A/(8\pi G_N)$ to $i\pi$ and pairing with its reflection assigns each state a relative weight $e^{-s}$.
Whenever the ordinary trace $Z=\tr(e^{-\hat{s}})$ is finite and nonzero on the selected space, the normalized density matrix $\rho$ and its von Neumann entropy satisfy
\begin{equation}
  \rho = Z^{-1}e^{-\hat{s}}, \qquad
  -\tr(\rho\log\rho) = \langle s\rangle_\rho + \log Z.
  \label{eq:entropy-id}
\end{equation}
The first term follows algebraically from the boost weight, whereas the value and large-area scaling of $\log Z$ depend on which microscopic states and multiplicities enter the trace.

\subsection{Local one-loop curvature coefficient and spin ensembles}

For a four-dimensional Schwarzschild black hole in pure Einstein gravity, gauge-fixed graviton and ghost fluctuations yield the local heat-kernel curvature coefficient $53/45$ in the convention with $(4\pi)^{-2}$ factored out~\cite{BastianelliBonezziMelis2022,Sen2013}.
Both Ref.~\cite{Chandrasekaran2026} (Section~5.3 and Appendix~D.1) and Sen's Euclidean analysis~\cite{Sen2013} use this same local curvature invariant.
At fixed reference units, converting it to the logarithmic correction of the density of states per unit mass summed over all angular momenta gives the large-area form
\begin{equation}
  S = s + \frac{106}{45}\log s + O(1).
  \label{eq:sen-entropy}
\end{equation}
Sen's ensemble analysis~\cite{Sen2013} also supplies calculable shifts when angular momentum is restricted: fixing one Cartan component at $J_3=0$ changes the $\log s$ coefficient to $167/90$, while restricting to rotational singlets ($J=0$) gives $77/90$ (\repo{audit/sen.md}{coefficient comparison}).

\subsection{Area-packet width and ensemble equivalence}

The Lorentzian trace proposed in Section~5.3 of Ref.~\cite{Chandrasekaran2026} incorporates the Gaussian area packet $f_b(s)$ of~\eqref{eq:reflected-weight}, whose standard deviation in $s$ is $\sqrt{2s_b}$.
Suppose the boost-weighted spectral density per unit $s$ behaves locally as $C s^p\,ds$ across an interior packet far from the endpoints $0$ and $s_*$.
Integrating over the packet gives
\begin{equation}
  \int C s^p \exp\!\left[-\frac{(s-s_b)^2}{4s_b}\right] ds
  \sim C\sqrt{4\pi}\, s_b^{p+1/2}.
  \label{eq:packet-integral}
\end{equation}
The packet width contributes a factor $\sqrt{s_b}$.
For four-dimensional Schwarzschild black holes, $s=4\pi G_N M^2/\hbar$ scales as mass squared, so converting a density per unit $s$ into a density per unit mass multiplies by the Jacobian $ds/dM \propto \sqrt{s}$---the exact same square-root factor.
Both counts therefore share a single logarithmic coefficient $p+\tfrac{1}{2}$ with no second half added; with $p+\tfrac{1}{2}=106/45$, the underlying boost-weighted per-unit-$s$ exponent is $p=167/90$ (numerically equal to Sen's fixed-$J_3$ per-unit-mass coefficient by a coincidence of four dimensions, though belonging to a different ensemble; see the \repo{audit/ensemble-crosscheck.md}{area-packet derivation}).

This match confirms the algebraic ensemble conversion of a shared local heat-kernel coefficient ($53/45$) rather than supplying an independent microscopic state count.
The conical surface terms for replica geometries are known for a specified entangling surface~\cite{FursaevPatrushevSolodukhin2013}. An independent Lorentzian derivation must still evaluate the full measure for the area collective coordinate, the constraints, and all ghost and zero-mode factors on joined replicas, determine whether the clock-filtered trace selects the marginal horizon saddle over a flat round sphere (where the surface integrand $R[q]+2\Theta_{(\ell)}\Theta_{(n)}$ integrates to $-8\pi$ rather than $+8\pi$), and establish the identification between the interior preparation and the asymptotic mass and spin ensemble.

\section{Operational Observables and Quantitative Benchmarks}
\label{sec:observables}

To anchor the theoretical discussion in specified preparations and calculable quantities, we summarize four concrete benchmarks (\repo{OBSERVABLES.md}{observable research program}) in Table~\ref{tab:observables} and the subsections below.
SI units are restored explicitly whenever macroscopic apparatus scales appear.

\begin{table}[htbp]
\centering
\caption{Summary of four operational benchmarks discussed in Section~\ref{sec:observables}, their numerical targets, and their relation to the vacuum four-dimensional proposal of Ref.~\cite{Chandrasekaran2026}.}
\label{tab:observables}
\small
\begin{tabular}{@{} >{\raggedright\arraybackslash}p{0.24\linewidth} >{\raggedright\arraybackslash}p{0.35\linewidth} >{\raggedright\arraybackslash}p{0.33\linewidth} @{}}
\toprule
Observable & Benchmark numerical target & Relation to Ref.~\cite{Chandrasekaran2026} \\
\midrule
Optical focusing vs.\ radar range & $|e|\simeq 2.34\times 10^{-26}$, $1-D/L\simeq 1.10\times 10^{-52}$, $L-D\simeq 4.39\times 10^{-49}\,\mathrm{m}$ & Classical Einstein benchmark; quantum deviation uncomputed \\
\addlinespace[4pt]
Gravitationally mediated spin entanglement & Entangling phase $-0.314\,\mathrm{rad}$, Bell overlap $0.578$ (separable ceiling $0.5$) & Coherent Newtonian matter coupling; 4D vacuum theory lacks matter \\
\addlinespace[4pt]
Sub-half-area Schwarzschild trace fraction & $2^{-257/90}\simeq 0.13816$ (packet) vs.\ $2^{-151/45}\simeq 0.09770$ (per-area) & Tests spectral measure once weights are derived independently \\
\addlinespace[4pt]
Boundary-frame boost interference & First zero at $\eta=0.5$ ($\Delta s=4\pi^2$ two-level) vs.\ $\eta=1$ (uniform band) & Distinguishes area distributions; blind to degeneracy at fixed area \\
\bottomrule
\end{tabular}
\end{table}

\paragraph{Optical focusing and radar range.}
Send a circular pencil of light counter-propagating against a vacuum plane gravitational wave and reflect its central ray back to the emitter's clock (\repo{audit/observable-optics.md}{optical derivation}).
On the wave axis, compare its round-trip radar range $L$ with the optical area distance $D=\sqrt{dA/d\Omega}$ inferred from the spot area $dA$ and calibrated launch solid angle $d\Omega$.
On a short constant-curvature segment, the signed spot ellipticity $e$ obeys
\begin{equation}
  1 - \frac{D}{L} = \frac{e^2}{5} + O(e^4).
  \label{eq:optics}
\end{equation}
For an illustrative strain $h=10^{-21}$ at frequency $f=100\,\mathrm{Hz}$ over $L=4\,\mathrm{km}$, Eq.~\eqref{eq:optics} yields $|e|\simeq 2.34\times 10^{-26}$ and a fractional optical-distance deficit $1-D/L\simeq 1.10\times 10^{-52}$.
The resulting distance change $L-D\simeq 4.39\times 10^{-49}\,\mathrm{m}$ is roughly $10^{31}$ times smaller than the differential arm-length response $hL$ of an optimally oriented interferometer at that strain, while the central round-trip echo time stays strictly unchanged in this geometry.
This is a classical general-relativistic benchmark for a screen and a clock; comparing quantizations requires computing a specific quantum deviation for the same apparatus.

\paragraph{Gravitationally mediated spin entanglement of falling masses.}
Prepare each of two $10^{-14}\,\mathrm{kg}$ masses independently on a pair of spin-labelled paths, with centres separated by $450\,\mu\mathrm{m}$ and a $250\,\mu\mathrm{m}$ splitting per mass; after a $2.5\,\mathrm{s}$ hold, recombine the paths and measure the two spins~\cite{Bose2017,MarlettoVedral2017} (\repo{audit/observable-lab.md}{spin observable and controls}).
Coherent Newtonian gravity gives an ideal entangling phase of $-0.314\,\mathrm{rad}$ and an overlap probability of $0.578$ with a specified entangled Bell state, exceeding the maximum $0.5$ achievable by any unentangled state.
A gravitational interpretation requires strict controls on electromagnetic forces, readout calibration, and initial correlations.
Predicting this signal from the four-dimensional vacuum construction of Ref.~\cite{Chandrasekaran2026} would first require coupling matter and specifying the detector model.

\paragraph{Sub-half-area trace fraction in the large Schwarzschild clock ensemble.}
In the proposed large Schwarzschild clock ensemble, form a ratio of boost-weighted physical traces with weight $e^{-\hat{s}}$, restricting the numerator to $s\le s_*/2$ and summing over all allowed areas $s\le s_*$ in the denominator (\repo{audit/observable-entropy.md}{trace ratio and assumptions}).
With the packet measure and large-diamond clock-ensemble scaling of Ref.~\cite{Chandrasekaran2026} (clock-ensemble exponent $257/90$, which belongs to the clock trace and is not directly comparable with Sen's fixed-mass numbers), the predicted fraction is $2^{-257/90}\simeq 0.13816$.
Assigning the same local one-loop power directly to a density per unit area instead gives $2^{-151/45}\simeq 0.09770$.
An independent test must derive the microscopic state weights and remove the regulator; an observational test would additionally require access to the prepared interior clock and area.

\paragraph{Interference of two boundary frames.}
If a reference apparatus could coherently boost one branch of a preparation by rapidity $\eta$, its squared fringe visibility would be $V(\eta)=|\langle e^{i\eta\hat{s}/(2\pi)}\rangle|^2$ (\repo{audit/observable-visibility.md}{visibility curves}).
For two equally weighted sharp areas separated by $\Delta s = 4\pi^2$ (about $158$ Planck areas), the first visibility zero occurs at $\eta=0.5$, whereas a uniform continuum of the same width gives its first zero at $\eta=1$.
These example preparations distinguish area distributions, but adding degenerate orthogonal states at each area while preserving its total probability leaves the fringes identically unchanged.
Moreover, the controlled-boost apparatus remains a theoretical proposal.

\section{Conjectures, Downstream Implications, and Open Problems}
\label{sec:open}

\subsection{Conditional implications of finite versus infinite rank}

The verified calculations of Sections~\ref{sec:classical}--\ref{sec:entropy} stand independently of the finite-rank conjecture in Section~5.2 of Ref.~\cite{Chandrasekaran2026}, whereas several headline consequences remain strictly conditional on proving a nonempty, finite-rank four-dimensional clock subspace (\repo{audit/dependencies.md}{dependency analysis}).

If a nonempty four-dimensional clock window has finite rank and total corner area preserves the selected subspace, then the restricted area operator has only finitely many distinct eigenvalues within that window.
Even so, their numerical spacing and degeneracies require separate calculations.
Extending the diamond construction to a finite thermal trace in anti-de Sitter (AdS) space likewise requires independent control of a closed thermal sum (unconditional convergence of the disconnected cylinder insertion in Section~5.2 of Ref.~\cite{Chandrasekaran2026} is equivalent to finiteness of the thermal trace).
Inferring spatial-topology redundancy or factorization from a two-boundary trace relies on additional completeness assumptions equating the physical norm of connected fillings with the Hilbert--Schmidt norm of separate boundary preparations.
Similarly, describing a black-hole evaporation history requires dynamical amplitudes for radiation escaping the diamond beyond an equilibrium two-sided correlator.

Conversely, an infinite-dimensional outcome for the causal diamond would still allow well-defined gravitational entropy calculations.
In algebraic approaches to semiclassical gravity, certain infinite-dimensional von Neumann observable algebras---such as type~II factors---admit a faithful semifinite trace $\tau$. A normal state with density $h$ has $\tau(h)=1$; the entropy $-\tau(h\log h)$, when finite, depends on the normalization of $\tau$ up to an additive constant~\cite{ChandrasekaranLongoPeningtonWitten2023,JensenSorceSperanza2023}.
Ref.~\cite{Chandrasekaran2026} instead works with a type~I factor and the ordinary trace (paralleling axiomatic Euclidean reconstructions~\cite{ColafranceschiDongMarolfWang2024}); obtaining a different algebraic type directly from the null-cone construction would require a separate derivation.

\subsection{Five open mathematical problems}

Settling the null-cone proposal turns on five concrete mathematical problems (\repo{REFEREE.md}{detailed technical questions}; \repo{audit/regulator-program.md}{regulator criteria}; \repo{audit/finite-kinematics.md}{physical corner-charge tests}):

\paragraph{Covariant closure and boundary flux of the finite-cap phase space.}
A complete classical formulation must specify the exact boundary data held fixed at both caps and the corner, re-derive the symplectic two-form from the full action including moving cap embeddings and joints, verify $d\omega=0$ on the endpoint variations that produce the $1/256$ rim defect of Section~\ref{sec:classical}, and verify that the Hamiltonian variation on the repaired phase space accounts for radiative boundary flux while recording which area and shape variations survive.

\paragraph{Continuum construction and ultraviolet control of the quantum clock.}
Starting with one physical polarization and one angular harmonic, the gravitational constraints must be solved prior to radial mode counting while keeping the cap fraction $\epsilon$, physical area, $G_N$, $\hbar$, and $\ell$ fixed, refining radial and angular cutoffs separately and then jointly.
Tracking both the bottom of the clock spectrum and the probability $\|P_\ell\Omega\|^2$ is decisive: a proof of a diverging spectral floor in the physical theory would rule out that bare clock prescription, whereas establishing nonzero limiting overlap requires controlling the quantum quadratic form and Hilbert-space embeddings across cutoffs (going beyond the $O(N^2)$ linearized mean squared shear and quadratic-model spectral floors of Section~\ref{sec:clock} and controlling high angular harmonics in interacting gravity).

\paragraph{Spectral rank in the physical space.}
Determining whether $P_\ell\HH$ has finite or infinite rank admits two complementary sufficient criteria.
Exhibiting a physical corner shape charge with a continuous spectral component on an allowed state whose spectral projections commute with $P_\ell$---incorporating edge excitations, cap constraints, and states reached by area-changing sources---would yield infinitely many orthogonal allowed states via the mechanism of Section~\ref{sec:clock}.
Conversely, constructing a finite-dimensional projection $F$ satisfying $\|(1-F)P_\ell\|<1$ ensures that $F$ is injective on $P_\ell\HH$, directly bounding $\rank(P_\ell)\le \rank(F)$.
A continuous classical parameter by itself leaves its quantum spectrum open.

\paragraph{Regulated interacting dynamics and independent closed sewing.}
Constructing regulated interacting evolution on a common state space requires proving that distinct intermediate cuts yield the same matrix element as the regulator is removed, and comparing the resulting weak-field four-graviton amplitude against the Einstein tree amplitude (and low-energy loop logarithms).
For closed contractions, the unnormalized boost response along distinct exhaustive cut orderings must be bounded uniformly against an independently defined uncut amplitude as more intermediate states are included.

\paragraph{Full Schwarzschild packet measure and saddle selection.}
Repeating the one-loop Schwarzschild determinant directly for the clock-filtered state requires checking whether the trace selects the marginal horizon saddle over a flat round sphere, evaluating the area-mode integration and zero modes on joined replicas under the same boundary conditions used in the state count, and testing whether the resulting Lorentzian trace recovers both the spin-summed coefficient $106/45$ and the spin-restricted values ($167/90$ and $77/90$) of Ref.~\cite{Sen2013}.

\paragraph{A two-copy test using finite sets of states.}
As a finite-cut diagnostic for the rank and closed-sewing problems above, let $F$ project onto $r$ orthonormal physical inputs inside the clock-selected subspace $P_\ell\HH$ at a fixed threshold, and suppose the reflected weight satisfies $0<mI\le W\le MI$ on $P_\ell\HH$ with fixed constants $m,M$ throughout refinement. Its finite Gram matrix $G_F=FWF|_{F\HH}$ is determined by open overlaps of fixed cap outputs and defines the conditioned two-copy return (purity) $Q_F=\tr(G_F^2)/(\tr G_F)^2$. For eigenvalues $\lambda\in[m,M]$, the inequality $(\lambda-m)(M-\lambda)\ge0$ and Cauchy--Schwarz give
\begin{equation}
  1\le rQ_F\le \frac{(m+M)^2}{4mM}.
  \label{eq:two-copy-bound}
\end{equation}
A uniform positive lower bound on $Q_F$ along an increasing family of projections exhausting $P_\ell\HH$ would therefore prove finite rank (while in infinite dimension $Q_F\to 0$ along every exhaustion). In an explicit continuous-area model, the regular replica integrals of Ref.~\cite{Chandrasekaran2026} (Appendix~D.3) give a return of $0.586222$, whereas for the same weight the microscopic return on any four or more orthonormal inputs is at most $0.400495$. Equating the replica integral to the microscopic return computed from physical states thus imposes an independently testable condition (\repo{audit/purity-test.md}{derivation and finite-matrix controls}).

\section{A Massless Scalar on the Two Null Sheets}
\label{sec:matter}

The four-dimensional state-count proposal of Ref.~\cite{Chandrasekaran2026} addresses \emph{vacuum} gravity; coupling a minimally coupled massless real scalar supplies three separate benchmarks---for exact classical geometry, semiclassical mean-metric backreaction, and conditional state counting---without deciding the vacuum conjecture itself.

\paragraph{Exact classical geometry.}
In spherical symmetry the gravitational shear vanishes. Along either null sheet, let the affine parameter $u$ increase from the tip toward the common waist of area radius $r_S$ with tangent $k=\partial_u$, and place the cap at $r_c=\epsilon r_S$ ($0<\epsilon<1$). Defining the cap-normalized clock reading $U=r'_c(u_S-u_c)$ with $r'_c\equiv r'(u_c)$ as in~\eqref{eq:affine-clock}, integrating the radial Einstein equation $r''=-4\pi G_N rT_{kk}$ twice gives the exact classical excess over a flat diamond of matched waist radius $r_S$,
\begin{equation}
  E = U - (1-\epsilon)r_S
    = 4\pi G_N\int_{u_c}^{u_S}(u_S-u)\,r(u)\,T_{kk}(u)\,du.
  \label{eq:matter-excess}
\end{equation}
For a classical scalar, $T_{kk}=(k^\mu\partial_\mu\phi)^2\ge 0$. In the spatially flat homogeneous Einstein--scalar solution with scale factor $a(t)\propto t^{1/3}$, choosing positive conformal tip times in the ratio $7:1$ (future to past) and cap fraction $\epsilon=5/16$ gives strictly positive expansion up to the waist on both sheets, with past- and future-tip readings $4025/5184$ and $5781/8000$ in units of $r_S$; both exceed the matched-waist flat value $1-\epsilon=11/16$ (\repo{audit/matter-diamond.md}{exact solution and weak-field limits}).

\paragraph{Semiclassical mean-metric response.}
By contrast, the renormalized stress expectation in a finite quantum state can take either sign along a finite null segment~\cite{FewsterRoman2003}. In a flat reference Minkowski diamond of waist radius $R$ (tips at $t=0,2R$), prepare the normalized spherical free-scalar mode $f(t,r)=R/(\pi[(R+it)^2+r^2])$. In units $c=\hbar=1$, the two states $(|0\rangle\pm 0.1|2_f\rangle)/\sqrt{1.01}$ share the same particle-number distribution and equal total mean energy $2/(101R)$, while their vacuum--pair interference enters the Minkowski vacuum-subtracted stress expectation $\langle{:}T_{kk}{:}\rangle$ with opposite signs. For $\epsilon=1/2$, substituting this smooth expectation value into the linearized semiclassical mean-metric version of~\eqref{eq:matter-excess} yields
\begin{equation}
  \frac{R E^{(1)}}{\ell_P^2} =
  \begin{array}{c|rr}
    & \text{minus state} & \text{plus state}\\ \hline
   \text{past-tip sheet} & -0.0114211 & +0.0159152\\
   \text{future-tip sheet} & -0.0035467 & +0.0071501
  \end{array}
  \qquad (R\gg\ell_P).
  \label{eq:scalar-sign}
\end{equation}
Here $E$ subtracts the flat reading at the \emph{physical} waist radius $r_S$, and its leading term $E^{(1)}$ samples $\langle{:}T_{kk}{:}\rangle$ against the kernel $(R-u)u$ on the reference sheet. This result tests a signed semiclassical mean-metric response (omitting graviton fluctuations); obtaining the expectation value or probability distribution of a full quantum clock requires a renormalized geometric operator and control of its fluctuations (\repo{audit/matter-null-energy.md}{packet normalization and integrals on both faces}).

\paragraph{Conditional scalar rank.}
The matter-coupled state count also depends on the allowed scalar boundary data. Suppose a nonempty physical clock projection $P$ is preserved by an exact noncompact scalar shift ($\phi\mapsto\phi+a$, $a>0$) and that a calibrated scalar-value measurement on $\mathbb{R}$ translates with it. If $P$ had finite rank $d$, the invariant state $P/d$ would induce a normalized probability measure on $\mathbb{R}$ invariant under a nonzero translation; its equal probabilities on the countably many translates of a fundamental interval cannot sum to unity, so finite rank is impossible under these hypotheses. Fixing the boundary value, compactifying the scalar target, breaking the symmetry, or removing the shift as a gauge redundancy alters the hypotheses and the physical count. Establishing the boundary algebra and joint symplectic form for the coupled quantum theory remains open, and none of these matter benchmarks by itself settles the vacuum four-dimensional conjecture (\repo{audit/matter-state-count.md}{scalar shift and finite-resolution test}).

\section{Boundary Repairs and Tests of the Continuum Limit}

The follow-up calculations below narrow specific classical boundary defects and isolate sharp operator criteria for the continuum limit, while underscoring why classical boundary recovery and fixed-background or fixed-resolution clock models remain strictly distinct from establishing a nonempty, finite-rank clock subspace in nonperturbative four-dimensional quantum gravity.

\paragraph{Variation of a moving rim.}
To match the notation of the boundary calculation, let $r$ denote the local corner area density relative to the chosen angular measure, $v=b$ the cap ratio, $m=rb^2$ the cap area density, and $z$ the two polarization coordinates with the target metric of Ref.~\cite{Chandrasekaran2026}, Eq.~(4.21). Stripping the factor $1/(16\pi G_N)$, the null shear potential on one generator has the form $\int a^2\langle\partial_a z,D_a z\rangle\,da$ in the absolute area radius $a=\sqrt r\,v$. Differentiating its physical endpoints, oriented from the common corner towards the cap, supplies the missing outer term
\begin{equation}
 \Delta\Omega_C=\frac b2\,Dm\wedge\langle z'(b),Dz_C\rangle .
\end{equation}
At fixed generator labels and fixed $b$, this term cancels the rim nonclosure of Section~\ref{sec:classical} for both polarizations. Removing that classical defect on the fixed-label radiative sector does not, however, complete the gravitational phase space or constrain the quantum state count: it still leaves the spacelike cap Cauchy data, their coupling to the rim, the endpoint inverse problem, and radiative flux across moving cuts to be incorporated. Even at fixed area, the variation of the clock at a cap with nonzero shear contains an endpoint evaluation that a Green operator inverting the bulk derivative on square-integrable interior sources cannot represent. The affine cut charge obeys a balance law involving the expansion and squared shear. Separately, the change of the cap symplectic form is governed by the symplectic current on the swept strip. Two vacuum radiative perturbations can have vanishing intrinsic jets at both end cuts while their integrated symplectic pairing on that strip is strictly positive (\repo{audit/rim-potential.md}{covariant rim variation}, \repo{audit/endpoint-inverse.md}{endpoint source}, \repo{audit/radiative-flux.md}{flux calculation}).

\subsection{A sharp clock on a joined cap and null sheet}
\label{sec:joint}
For cuts $a<b<c$ ordered by $v=\rho/\rho_S$, let $p_b=(d\rho/du)_b$ in one affine parameter $u$. The readings $T_{bc}=p_b(u_c-u_b)$ obey $T_{ac}=T_{ab}+(p_a/p_b)T_{bc}$.
Calibrating both segments at $a$ makes the focusing kernel continuous across $b$, so internal shear evaluations cancel upon integration by parts against smooth polarization variations.
Resetting the affine calibration at the cap rim $C$ ($v=b$, with waist $S$ at $v=c=1$), as in Ref.~\cite{Chandrasekaran2026}, Eq.~(3.27), introduces a jump at $C$ and produces the rim differential $-\rho_C(u_S-u_C)\sigma^{AB}_{(\ell)}\delta q^{\mathrm{TF}}_{AB}/2$ at fixed local area density and generator labels, where $\sigma^{AB}_{(\ell)}$ is determined by the tangential trace-free cap momentum and the second fundamental form of $C$ within the spacelike cap (\repo{audit/joint-edges.md}{cap matching}).
The inverse problem asks for a single Hamiltonian tangent whose symplectic pairing reproduces the differential against every admissible test (\repo{audit/joint-clock.md}{joined bracket and scope}).

Consider the vacuum plane wave $ds^2=-2\,dU\,dV+A_1(U)^2dx_1^2+A_2(U)^2dx_2^2$ with $A_1=\cosh U$ and $A_2=\cos U$, placing both smooth tips on the central geodesic at $(U,V)=(0,0)$ and $(1,1)$ ($L=D=1$) and taking cap fraction $\epsilon=1/2$, while tracking the motion of both the waist intersection and the cap.
A compactly supported smooth pulse of shrinking width $h\downarrow 0$ in $H=\log(A_1/A_2)$ with zero total integral and unit integral over its earlier half defines a tangent to an exact vacuum family after solving the focusing equation for $R=\sqrt{A_1A_2}$.
Near the axial direction, where $U$ has a nondegenerate maximum on the smooth cap, angular integration samples only the earlier half of the pulse, yielding the limiting first variation $-0.060534$ in these units.
On a spacelike ball spanning the waist, $U$ is a smooth transverse coordinate through the pulse, also on the boundary of the ball. Both halves enter every smooth metric symplectic pairing there, which tends to zero.
Conservation transports that integral to the joined cap and null sheet whenever the metric and boundary frame are shared across the seam, ruling out a smooth metric Hamiltonian vector for $\bar u_*$ at this curved background under the shared-frame symplectic form (\repo{audit/joint-focal.md}{exact limit and numerical check}).
For a smooth relative displacement of the two embeddings across this curved rim, the limiting covariant charge is proportional to the difference of the two fractional principal stretches of the induced metric at the axial point; matching both induced metric components removes this contribution, and the cap area--boost term also tends to zero for a spacelike collar with prescribed tilt there (\repo{audit/match-curved.md}{curved-rim charge calculation}).
A distributional inverse can act on smooth probes. In a polarized Einstein--Rosen sector (toroidal cuts, fixed areal embeddings), a smooth physical disturbance shifts the clock by $-3/2$ entirely through its rim term, and the causal Peierls inverse vanishes in the open cap while carrying a delta distribution in the initial canonical momentum at the rim (\repo{audit/joint-peierls.md}{distributional inverse}).

\paragraph{One-sided charge, common smooth seam, and separately marked relative angle.}
Let $M=\langle\rho_S^2\rangle/\langle\rho_S\rangle$, with brackets denoting the mean over the emitter's sky, and let $D=\bar u_*-(1-a)M$ be the clock excess for a cap at $v=a$. Focusing is quadratic in shear, so $D$ is critical ($dD=0$) on a round flat diamond; adjusting the future tip to hold $M=R$ makes $D$ differ from the raw clock $\bar u_*$ by a constant. Work in units $R=1$ with $a=1/3$, and consider two smooth transverse-traceless vacuum quadrupoles sharing a screen tensor $B$ with $\langle B:B\rangle=3/5$. For a fixed reference wave $W$ with unit radial strain slope at $a$ and a single narrowing packet $Z_\eta$ with unit strain at $a$, integration by parts gives the single-packet mixed second variation
\begin{equation}
 \lim_{\eta\to0} d^2D(W,Z_\eta)
 =-\tfrac14a(1-a)\langle B:B\rangle=-\tfrac1{30}.
 \label{eq:joint-hessian}
\end{equation}
The packet pairs to zero in the limit against any smooth Einstein metric tangent on the waist disk. By contrast, a \emph{one-sided} sheet endpoint with an independent angular displacement $V=G_N\nabla P_2(\cos\theta)$ can supply this same $-1/30$ through its normal-connection charge (\repo{audit/joint-flat-hessian.md}{vacuum wave and clock}, \repo{audit/joint-angular.md}{one-sided charge}).

On a \emph{common smooth seam}, however, a shared displacement carries zero net charge. Let $T$ denote the future cap normal, $s$ its outward spatial normal, and $\ell=e^\beta(T+s)$ the outgoing null normal, so that the null normal connection obeys $\varpi_A=K_{As}+\partial_A\beta$ with $K_{As}=g(\nabla_A T,s)$. With positive sheet-end orientation, the sheet, cap, and joint angular densities cancel pointwise: $\mu_C(\varpi_A-K_{As}-\partial_A\beta)/(8\pi G_N)=0$, provided both faces refer to the same tangential displacement and normal boost derivative. A change of cap slicing with lapse vanishing at the rim likewise has cancelling cap and joint pairings on the vacuum constraints (\repo{audit/match-action.md}{covariant charge}, \repo{audit/match-boost.md}{cap tilt}). By contrast, for two \emph{independently marked} faces, let a comparison map $Y$ identify the rims with $q_+=Y^*q_-$ and $\ell_+=e^\gamma Y^*\ell_-$. Normal-frame transport and the variation of metric matching give
\begin{equation}
 \varpi_+=Y^*\varpi_-+d\gamma,\qquad
 \delta q_+-Y^*\delta q_-=\mathcal L_w q_+,\qquad
 w=(dY)^{-1}\delta Y ,
 \label{eq:marked-cap-matching}
\end{equation}
where $\delta q_\pm$ are the metric variations in the two side charts. Holding those marked variations equal ($\delta q_+=Y^*\delta q_-$) restricts $w$ to a Killing field of $q_+$. Allowing relative normal motion on the round rim adds only a pure trace, so the trace-free quadrupolar stretch of $V$ above still cannot survive when the marked variations have equal trace-free parts. A pair of independently anchored frames can retain angular momentum along an isometry even with fixed marked metrics. The non-Killing quadrupolar slide requires marked variations whose difference supplies its strain (\repo{audit/match-reduction.md}{comparison-map variation}).

Kerr exterior spheres on a Boyer--Lindquist time slice~\cite{BoyerLindquist1967} cleanly separate these cases. Reversing the spin preserves the marked two-metric on the sphere, while in the cap-adapted frame the one-sided charge of the axial Killing field is $(8\pi G_N)^{-1}\int_C\mu_C\varpi_\phi=-J$, where $J$ is the signed angular momentum. Across a smooth seam in that same Kerr geometry, the adjacent face contributes $+J$, so the common rotation has zero net charge; joining opposite spins by the identity marking would violate angular momentum matching. Because a regular axisymmetric vacuum ball with no inner boundary has zero axial flux by the momentum constraint, this Kerr example specifically tests a vacuum exterior collar with an inner boundary (\repo{audit/match-rotation.md}{metric and exact charge integral}).

\paragraph{Two packets with unchanged local frame data.}
Take the flat tips at $(t,\mathbf x)=(0,\mathbf0),(2,\mathbf0)$. A vacuum quadrupole generated from the regular radial wave $[F(t+r)-F(t-r)]/r$ induces first-order conformal screen strain $G_F(v)B$ on the outgoing cone, with
\begin{equation}
 G_F(v)=\frac{F^{(4)}(2v)}{v}-\frac{2F^{(3)}(2v)}{v^2}+\frac{F''(2v)}{v^3},
 \qquad d^2D(W,Z)=\frac3{20}\int_a^1 v(1-v)G'_W G'_Z\,dv ,
 \label{eq:packet-hessian}
\end{equation}
where $F^{(k)}(2v)$ differentiates $F(u)$ with respect to its argument $u=2v$ and primes on $G_W,G_Z$ denote $v$-derivatives. For the test sequence $Z_\eta$, choose $F(2v)$ to be the difference (outer minus inner) of two smooth bells of unit integral in $v$, each of half-width $\eta$, centred at $a+2\eta$ and $a-2\eta$. Both profiles vanish on an open gap around the rim $v=a$. On that gap, optical transport returns exactly to its flat value at first order: the trace-free optical strain vanishes, the normal-connection transport has primitive $\int_0^v tG_F(t)\,dt=F^{(3)}(2v)/2-F''(2v)/v$, and linearized vacuum focusing fixes the area at fixed affine parameter. Consequently, the first variations of the cut embeddings and of the transported null tetrads vanish on collars of both rims and the waist, so independent local frame marks may be held fixed along $Z_\eta$ (\repo{audit/match-packets.md}{ray and tetrad transport}).

To isolate a mixed second derivative distinct from the single-packet value $-1/30$ in~\eqref{eq:joint-hessian}, choose a fixed reference wave $W$ with $G'_W(v)=(v-a)^5+O((v-a)^6)$, generated from a master profile $F(u)$ proportional to $(u-2a)^{10}$ near $u=2a$. Its rim strain and slope vanish through fourth order, so the single-packet angular term does not enter. Integrating~\eqref{eq:packet-hessian} by parts selects only the outer bell (since the clock integrates solely over $v\in[a,1]$), yielding the distinct limiting mixed second derivative
\begin{equation}
 \lim_{\eta\to0} d^2D(W,Z_\eta)
 =-\frac{3}{20}\frac{(1-a)5!}{16}=-\frac34 .
 \label{eq:compensated-hessian}
\end{equation}
By contrast, on the spacelike waist disk $t=1$, both bells land in an interior annulus near $r=1/3$ with opposite signs; their signed sum and all its derivatives converge to zero as distributions as $\eta\to0$, driving the bulk Einstein symplectic pairing with every fixed smooth metric tangent to zero. When both faces share one smooth metric across the rim, conservation~\cite{IyerWald1994} transports this vanishing pairing to the joined cap and null sheet. The argument only requires that the derivative of the proposed generator along the fixed reference wave be a smooth tangent. The tests can grow in energy as they narrow, since the Hamilton equation must hold for each smooth test separately. Smooth embedding variations in the extended potential (Ref.~\cite{Chandrasekaran2026}, Eq.~(B.21)) contribute on shell only through cut terms, and every local cut term of finite differential order pairs to zero on $Z_\eta$ once independent local marks are held fixed on the unchanged collars. On a regular reduction, a $C^1$ family of smooth Hamiltonian vectors $X_D$ must vanish at the critical flat solution, so differentiating its Hamilton equation along $W$ requires $(d_W X_D)_{\rm flat}$ to reproduce~\eqref{eq:compensated-hessian} through just such a bulk-plus-local pairing. The nonzero limit $-3/4$ therefore excludes a $C^1$ smooth assignment, although the zero vector still solves the Hamilton equation at flat space itself.

\paragraph{Finite collections of boundary responses.}
A nonlocal frame transported through the interior of the spacelike cap can still detect the inner packet even when all local fields are unchanged on a rim collar. To test such frames, let $f(v)=F(2v)$ range over all smooth compactly supported master profiles on an interval about $a$, and suppose the additional boundary two-form on $(X,Z_f)$ factors through a fixed finite list of linear functionals of $f$ for all smooth responses $X$ at flat space. This covers any two-form pulled back from finitely many tensor fields on the round cuts whose linearized response is rotation-equivariant at the flat background, since each such field carries only finitely many components in the quadrupole representation. Taking reference waves $W_n$ with $G'_{W_n}(v)=(v-a)^{n+5}+O((v-a)^{n+6})$ for $n\ge0$ and integrating~\eqref{eq:packet-hessian} by parts yields, against $f(v)\,dv$, the response density
\begin{equation}
 \mathbf1_{v>a}\big[c_n(v-a)^n+O((v-a)^{n+1})\big],\qquad
 c_n=-\frac{3(1-a)}{320}\frac{(n+5)!}{n!}.
 \label{eq:cap-hessian-rank}
\end{equation}
In any nonzero finite linear combination, the lowest participating order $n$ exhibits a jump in its $n$-th derivative at $v=a$, so the densities~\eqref{eq:cap-hessian-rank} are linearly independent modulo smooth densities. Because the waist-disk Einstein pairing supplies a smooth radial density for every smooth metric tangent while the extra boundary two-form spans at most the assumed finite-dimensional space, together they cannot represent all derivatives $(d_{W_n}X_D)_{\rm flat}$ of a $C^1$ smooth Hamiltonian field. A reduction restricting the admissible waves or introducing infinitely many independent radial boundary responses requires a separate range calculation (\repo{audit/match-packets.md}{rank proof and coefficient checks}).

\paragraph{A single clock with freely varying waist.}
Holding $M$ fixed by adjusting the later tip defines a restricted preparation; inferring an ambient Hamiltonian vector from a restriction requires a separate tangency argument. Instead, set $r_A=\sqrt{A/(4\pi)}$. At a round waist $\rho_S=R$, one has $dM=dr_A$ and
\begin{equation}
 d^2(M-r_A)(W,Z)=\frac{\langle\delta_W\rho_S\,\delta_Z\rho_S\rangle
       -\langle\delta_W\rho_S\rangle\langle\delta_Z\rho_S\rangle}{R}.
\end{equation}
At a round waist $dM=dr_A$ makes the area-subtracted clock critical on every tangent. All quadrupole test packets supported near the rim satisfy $\delta_Z\rho_S=0$ pointwise on the waist, so the mixed Hessians of $\bar u_*-(1-a)r_A$ and $D$ coincide for these tests. If the exterior waist retains the smooth uniform frame boost with Hamiltonian $A/(8\pi G_N)$, then $r_A$ likewise admits a smooth Hamiltonian vector; a $C^1$ smooth field for $\bar u_*$ on this freely varying phase space would therefore induce one for the critical difference $\bar u_*-(1-a)r_A$, contradicting~\eqref{eq:compensated-hessian} or~\eqref{eq:cap-hessian-rank} under their respective hypotheses. This step relies only on the total-area boost of Ref.~\cite{Chandrasekaran2026}, Eqs.~(3.50b)--(3.51); gluing away the exterior frame or tying it to a cap-fixed affine frame alters that hypothesis (\repo{audit/match-area.md}{charge and free-waist proof}).

A smooth source allows a direct Peierls calculation. In a globally hyperbolic vacuum Kasner spacetime with compact spatial $T^3$ sections, a smooth gauge-invariant probe of the torus shape has a nonzero classical Peierls bracket with the matched-radius clock \emph{averaged over emitter positions on the spatial torus}, with leading coefficient $-0.00036096\,R^5/T^4$ at $\epsilon=1/2$ (launch time $T$).
At this symmetric background, the spatial position average and torus symmetries also collapse the clock differential to a smooth Hamiltonian tangent---locally a spatial rescaling invisible to gauge-invariant probes in a contractible patch and detected only by the global torus modulus (\repo{audit/joint-jacobi.md}{source and Hamiltonian response}).
For the single sharp clock, a distributional inverse on a specified domain, its quantization and the rank of its spectral projector remain open. From the repository root, \texttt{python3 -B checks/match\_packets.py} reproduces the limit $-3/4$ and the first two higher-order coefficients $-9/2$ and $-63/4$.

\subsection{Existence and rank of two specified clocks}
Two continuum model clocks give opposite, rigorous answers within their own assumptions, illustrating why neither a frozen-background obstruction nor a fixed-resolution construction decides finite rank for the full four-dimensional theory. First, in a bare model with the null bracket frozen on a smooth curved polarization background, one retains both polarizations, the nonlinear target metric in the focusing, and the outer exponential; the weak uncertainty bound for two radial currents then yields a state-independent focusing floor proportional to the square of the radial cutoff, emptying every fixed finite-reading window. Second, a continuum construction in linearized gravity smooths the shear in angle and radius at a fixed physical resolution before squaring it; Wick contraction bounds the vacuum exponential moment, yielding a vacuum acceptance lower bound of $0.6373$ at $s=10{,}000$, cap ratio $1/4$, and a ceiling $1.2$ times the shear-free reading, while exciting sufficiently weakly coupled modes proves infinite rank above the vacuum mean. Furthermore, restoring the nonlinear hyperbolic target metric introduces genuine domain instability: quantizing a smooth fixed-density radiative truncation gives a canonical pair through a late pulse, while an arbitrarily small second polarization on the earlier patches collapses the dense finite-form domain to $\{0\}$. Transferring that result requires controlling quantization before restriction. Together, these examples isolate a dynamical question about the ultraviolet response of the clock in Ref.~\cite{Chandrasekaran2026}: raising precision at fixed smoothing width and removing that width are distinct limits, and controlling the nonperturbative gravitational constraints and operator domains remains open (\repo{audit/clock-spectrum.md}{bare model}, \repo{audit/clock-construction.md}{smoothed model}, \repo{audit/clock-domain-jump.md}{nonlinear domain counterexample}).

\paragraph{A sufficient test for the state count.}
Two complementary operator criteria sharpen the four-dimensional rank problem without presupposing closed sewing. On the finite-rank side, for a unit vector $\psi$ and the full irreducible source algebra, the closure of $\{PA\psi:\|A\|\le1\}$ is the closed unit ball of the clock projection $P$; uniform approximation of this unit ball by a single finite-dimensional subspace with error strictly less than one forces $P$ to have finite rank, whereas checking only a restricted cyclic subalgebra or a dense collection of preparations without the unit-ball bound can yield false positives. On the infinite-rank side, a covariantly transformed marked corner area moment has no eigenvectors and forces infinite rank if it commutes with a nonzero clock window, though nonlinear canonical ordering can break classical boost-invariance commutation. Establishing either dynamical estimate on the full gravitational source algebra remains unresolved (\repo{audit/sewing-repair.md}{finite-rank criteria}, \repo{audit/corner-commutant.md}{independent area-charge obstruction}).

\paragraph{A classical interaction and a horizon-sector control.}
Two counterpropagating, aligned plane waves in vacuum furnish an exact classical self-interaction benchmark linear in a weak probe of strength $\delta$ and all orders in the pump before focusing. For incoming log strains $\varepsilon f(u)$ and $\delta g(v)$, the local transverse log-length ratio $V$ satisfies $\partial_\delta V|_0=g(v)/r_f(u)$, where $r_f''=-\varepsilon^2 f'^2 r_f/4$ with initial radius one and zero derivative. For $f(u)=\sin^4(\pi u)$ on $[0,1]$, the gain at pump exit is $1+(5\pi^2/64)\varepsilon^2+O(\varepsilon^4)$. The result fixes both the nonlinear observable and the affine frame, providing a classical finite-region target distinct from quantum multi-graviton scattering (\repo{audit/interacting-test.md}{wave transfer and tidal response}). In exact Schwarzschild spacetime, marginality uniquely selects the horizon sphere among radial graphs on a fixed ingoing null sheet; cuts along the horizon itself can nevertheless share the same intrinsic round area while carrying different inward shear and affine clock readings. Selecting the marginal saddle, its clock-filtered trace weight, and its collective-coordinate measure thus remains a separate calculation (\repo{audit/entropy-determinant.md}{shape determinant and deformed horizon cuts}).

\section{Conclusions}
\label{sec:conclusions}

Our analysis separates what is established in Ref.~\cite{Chandrasekaran2026} from the defects and proof gaps that leave its central finite-dimensionality claim open.
On the invertible-tetrad branch, the constrained first-order bulk equations reproduce classical vacuum Einstein dynamics, and null shear focuses rays by the exact affine delay~\eqref{eq:affine-clock} and area bound~\eqref{eq:clock-bound}.
The two-dimensional JT gravity control confirms that a bounded clock alone permits infinitely many orthogonal states.
In four dimensions, the specified single-harmonic Gaussian pair overlap~\eqref{eq:pair-sum} sums exactly to~\eqref{eq:pair-limit} and exhibits a genuine nonuniform weak-coupling expansion~\eqref{eq:pair-asymp}.
Around a regular classical saddle, the stationary-phase equations and fixed-resolution one-loop response have the expected form (with continuum renormalization left open), and the local graviton-and-ghost curvature coefficient $53/45$ converts arithmetically to Sen's spin-summed logarithmic coefficient $106/45$.

The finite-cap defect calls for a boundary two-form derived with the cap and its rim allowed to vary. Independently, the compensated vacuum-packet family excludes a $C^1$ assignment of smooth Hamiltonian vectors through flat space for the clock excess under the matching and boundary hypotheses of Section~\ref{sec:joint}, and retaining the exterior total-area boost extends this obstruction to a single sharp clock with freely varying waist. The spectral floors found in two regulated shear models motivate a separate continuum test of the gravitational clock. Even a construction of a nonzero clock projection would leave a second problem: proving a finite closed trace independently of the proposed state count. Establishing that trace and deriving the area-packet measure would permit a genuinely microscopic comparison with the Schwarzschild entropy.

\section{Reproducibility, Data, and Authorship Disclosure}
\label{sec:repro}

All companion derivations, the \repo{EQUATIONS.md}{equation index}, the \repo{audit/verification-scope.md}{verification inventory by section}, the \repo{audit/dependencies.md}{dependency analysis}, the \repo{reviews/README.md}{public review record}, and the \repo{docs/paper.tex}{LaTeX source} of this article are archived in the public repository at \url{https://github.com/sprajs/quantum-gravity}.
The four core numerical checks cited in Sections~\ref{sec:classical}--\ref{sec:sewing} rely solely on the Python~3 standard library and are listed in Table~\ref{tab:repro}; running \texttt{make check} from the repository root executes the full verification suite (the Appendix B checks invoke SymPy via \texttt{uv}).

\begin{table}[htbp]
\centering
\caption{Executable Python~3 standard-library scripts reproducing the key numerical checks in Sections~\ref{sec:classical}--\ref{sec:sewing}.}
\label{tab:repro}
\small
\begin{tabular}{@{} l >{\raggedright\arraybackslash}p{0.58\linewidth} @{}}
\toprule
Command & Verified quantitative output \\
\midrule
\texttt{python3 checks/03-3.py} & Finite-cap exterior derivative on the polynomial test family vanishes at $b=0$ and equals $1/256$ at $b=1/2$ in exact rational arithmetic. \\
\addlinespace[4pt]
\texttt{python3 checks/pair.py} & Direct mode sum through $N=100{,}000$ matches independent evaluation of~\eqref{eq:pair-limit} within analytic tail bounds ($|C_N-C|\simeq 2.50\times 10^{-12}$ at $x=0.1$). \\
\addlinespace[4pt]
\texttt{python3 checks/corner\_clock.py} & At $\ell/u_0=1.1$, accepted eigenvalue count grows from $6$ ($n=16$) to $94$ ($n=4{,}096$) across truncated matrix sizes (Table~\ref{tab:corner-clock}). \\
\addlinespace[4pt]
\texttt{python3 checks/05-2.py} & Three exhaustive cut orderings in the two-area model yield boost contrasts $-0.45970$, $-0.18388$, and $-0.02242$. \\
\bottomrule
\end{tabular}
\end{table}

\paragraph{AI authorship and independence disclosure.}
This study analyzes the 83-page version~1 of Chandrasekaran's preprint~\cite{Chandrasekaran2026}, posted on 25~November~2026.
Both the assessed manuscript and the present study acknowledge AI assistance.
In its acknowledgements, Ref.~\cite{Chandrasekaran2026} states that several of its calculations and proofs were developed through collaboration with GPT-6 Astra.
In this study, Astra led the ongoing calculations, technical cross-checks, and drafting; Fable provided independent technical reviews of the earlier audit from a separate model family. The subsequent gap calculations and their technical rechecks were performed by Astra; and Neon completed public editorial revisions for this edition.
Because Astra checks share a model family with the preprint's collaborator, and because this study has not had human specialist peer review, readers should evaluate every claim directly from its explicit hypotheses, derivations, counterexamples, and executable scripts against the versioned record~\cite{Chandrasekaran2026}.

\begingroup
\small
\raggedright
\begin{thebibliography}{20}\setlength{\itemsep}{2pt plus 1pt}

\bibitem{Chandrasekaran2026}
V.~Chandrasekaran,
\emph{An algebraic quantization of gravity from the null cone},
\href{https://arxiv.org/abs/2611.34924v1}{arXiv:2611.34924v1} (25~November~2026), 83~pp.

\bibitem{Raychaudhuri1955}
A.~Raychaudhuri,
\emph{Relativistic cosmology. I},
Phys.\ Rev.\ \textbf{98}, 1123 (1955),
\href{https://doi.org/10.1103/PhysRev.98.1123}{doi:10.1103/PhysRev.98.1123}.

\bibitem{Sachs1961}
R.~K.~Sachs,
\emph{Gravitational waves in general relativity. VI. The outgoing radiation condition},
Proc.\ R.\ Soc.\ Lond.\ A \textbf{264}, 309--338 (1961),
\href{https://doi.org/10.1098/rspa.1961.0202}{doi:10.1098/rspa.1961.0202}.

\bibitem{Reisenberger2008}
M.~P.~Reisenberger,
\emph{The Poisson bracket on free null initial data for gravity},
Phys.\ Rev.\ Lett.\ \textbf{101}, 211101 (2008),
\href{https://arxiv.org/abs/0712.2541}{arXiv:0712.2541}.

\bibitem{Reisenberger2013}
M.~P.~Reisenberger,
\emph{The symplectic 2-form for gravity in terms of free null initial data},
Class.\ Quantum Grav.\ \textbf{30}, 155022 (2013),
\href{https://arxiv.org/abs/1211.3880}{arXiv:1211.3880}.

\bibitem{ChandrasekaranFlanagan2023}
V.~Chandrasekaran and \'E.~\'E.~Flanagan,
\emph{Horizon phase spaces in general relativity},
\href{https://arxiv.org/abs/2309.03871}{arXiv:2309.03871} (2023).

\bibitem{BoyerLindquist1967}
R.~H.~Boyer and R.~W.~Lindquist,
\emph{Maximal analytic extension of the Kerr metric},
J.\ Math.\ Phys.\ \textbf{8}, 265--281 (1967),
\href{https://doi.org/10.1063/1.1705193}{doi:10.1063/1.1705193}.

\bibitem{IyerWald1994}
V.~Iyer and R.~M.~Wald,
\emph{Some properties of Noether charge and a proposal for dynamical black hole entropy},
Phys.\ Rev.\ D \textbf{50}, 846--864 (1994),
\href{https://arxiv.org/abs/gr-qc/9403028v1}{arXiv:gr-qc/9403028v1}.

\bibitem{GoroffSagnotti1986}
M.~H.~Goroff and A.~Sagnotti,
\emph{The ultraviolet behavior of Einstein gravity},
Nucl.\ Phys.\ B \textbf{266}, 709--736 (1986),
\href{https://doi.org/10.1016/0550-3213(86)90193-8}{doi:10.1016/0550-3213(86)90193-8}.

\bibitem{DunbarJehuPerkins2017}
D.~C.~Dunbar, G.~R.~Jehu, and W.~B.~Perkins,
\emph{Two-loop gravity amplitudes from four dimensional unitarity},
Phys.\ Rev.\ D \textbf{95}, 046012 (2017),
\href{https://arxiv.org/abs/1701.02934}{arXiv:1701.02934}.

\bibitem{Behrndt2010}
J.~Behrndt, S.~Hassi, H.~de~Snoo, and R.~Wietsma,
\emph{Monotone convergence theorems for semi-bounded operators and forms with applications},
Proc.\ Roy.\ Soc.\ Edinburgh Sect.\ A \textbf{140}, 927--951 (2010),
\href{https://doi.org/10.1017/S030821050900078X}{doi:10.1017/S030821050900078X}.

\bibitem{Bekenstein1973}
J.~D.~Bekenstein,
\emph{Black holes and entropy},
Phys.\ Rev.\ D \textbf{7}, 2333--2346 (1973),
\href{https://doi.org/10.1103/PhysRevD.7.2333}{doi:10.1103/PhysRevD.7.2333}.

\bibitem{Hawking1975}
S.~W.~Hawking,
\emph{Particle creation by black holes},
Commun.\ Math.\ Phys.\ \textbf{43}, 199--220 (1975),
\href{https://doi.org/10.1007/BF02345020}{doi:10.1007/BF02345020}.

\bibitem{BastianelliBonezziMelis2022}
F.~Bastianelli, R.~Bonezzi, and M.~Melis,
\emph{Gauge-invariant coefficients in perturbative quantum gravity},
Eur.\ Phys.\ J.\ C \textbf{82}, 1139 (2022),
\href{https://arxiv.org/abs/2206.13287}{arXiv:2206.13287}.

\bibitem{FursaevPatrushevSolodukhin2013}
D.~V.~Fursaev, A.~Patrushev, and S.~N.~Solodukhin,
\emph{Distributional geometry of squashed cones},
Phys.\ Rev.\ D \textbf{88}, 044054 (2013),
\href{https://arxiv.org/abs/1306.4000}{arXiv:1306.4000}.

\bibitem{Sen2013}
A.~Sen,
\emph{Logarithmic corrections to Schwarzschild and other non-extremal black hole entropy in different dimensions},
JHEP \textbf{04} (2013) 156,
\href{https://arxiv.org/abs/1205.0971}{arXiv:1205.0971}.

\bibitem{Bose2017}
S.~Bose \emph{et al.},
\emph{Spin entanglement witness for quantum gravity},
Phys.\ Rev.\ Lett.\ \textbf{119}, 240401 (2017),
\href{https://arxiv.org/abs/1707.06050}{arXiv:1707.06050}.

\bibitem{MarlettoVedral2017}
C.~Marletto and V.~Vedral,
\emph{Gravitationally induced entanglement between two massive particles is sufficient evidence of quantum effects in gravity},
Phys.\ Rev.\ Lett.\ \textbf{119}, 240402 (2017),
\href{https://arxiv.org/abs/1707.06036}{arXiv:1707.06036}.

\bibitem{ChandrasekaranLongoPeningtonWitten2023}
V.~Chandrasekaran, R.~Longo, G.~Penington, and E.~Witten,
\emph{An algebra of observables for de Sitter space},
JHEP \textbf{02} (2023) 082,
\href{https://arxiv.org/abs/2206.10780}{arXiv:2206.10780}.

\bibitem{JensenSorceSperanza2023}
K.~Jensen, J.~Sorce, and A.~J.~Speranza,
\emph{Generalized entropy for general subregions in quantum gravity},
JHEP \textbf{12} (2023) 020,
\href{https://arxiv.org/abs/2306.01837}{arXiv:2306.01837}.

\bibitem{ColafranceschiDongMarolfWang2024}
E.~Colafranceschi, X.~Dong, D.~Marolf, and Z.~Wang,
\emph{Algebras and Hilbert spaces from gravitational path integrals: Understanding Ryu--Takayanagi/HRT as entropy without AdS/CFT},
JHEP \textbf{10} (2024) 063,
\href{https://arxiv.org/abs/2310.02189}{arXiv:2310.02189}.

\bibitem{FewsterRoman2003}
C.~J.~Fewster and T.~A.~Roman,
\emph{Null energy conditions in quantum field theory},
Phys.\ Rev.\ D \textbf{67}, 044003 (2003),
\href{https://doi.org/10.1103/PhysRevD.67.044003}{doi:10.1103/PhysRevD.67.044003}.

\end{thebibliography}
\endgroup

\end{document}
